\section{Marco teórico y trabajos relacionados}

La programación orientada a objetos ha sido uno de los pilares fundamentales en la transformación del desarrollo de software \cite{oop2023principles}. Su enfoque en la encapsulación, la herencia y el polimorfismo ha permitido construir modelos de software más coherentes y mantenibles. Ejemplo de ello es la implementación computacional del método de red de vórtices inestacionario, desarrollada mediante paradigmas de orientación a objetos y co-simulación \cite{oop2022simulation}. Este tipo de enfoques demuestra cómo la modularidad y la abstracción permiten desarrollar sistemas complejos que integran simulaciones físicas, procesamiento matemático y visualización, sin comprometer la claridad del código ni la independencia de los módulos.

Los conectores UML 2.0 han permitido una descripción más detallada de los aspectos estructurales y conductuales de los sistemas \cite{uml2020guide}, haciendo posible representar las interacciones entre componentes, servicios y flujos de datos de manera estandarizada \cite{uml2021agile}. Esta capacidad de modelado es fundamental para la documentación y comunicación de arquitecturas complejas en equipos de desarrollo.

En el ámbito metodológico, la integración de la metodología DAC (Development Agile Cycle) con los estándares PMBOK \cite{pmbok2017} y CMMI-DEV \cite{cmmi2018} ha marcado una pauta importante en la gestión de proyectos de software \cite{dac2020methodology}. Este modelo mixto reconoce la necesidad de flexibilidad del ágil, pero la combina con la estructura y trazabilidad del enfoque tradicional. En entornos de desarrollo con múltiples equipos o proyectos a gran escala, esta sinergia garantiza que los productos no solo se entreguen a tiempo, sino que cumplan con métricas de calidad, madurez y documentación exigidas por la industria.

Los patrones de diseño y las arquitecturas reutilizables representan otra pieza esencial en este entramado \cite{gamma1995patterns}. Entre ellos, el patrón Modelo-Vista-Controlador (MVC) se ha consolidado como uno de los más influyentes \cite{mvc2019analysis}, particularmente en frameworks como Laravel y React. MVC separa la lógica de negocio, la gestión de datos y la interfaz de usuario, promoviendo una arquitectura limpia, mantenible y escalable.

Dentro de la misma línea de patrones, el patrón Factory ha cobrado relevancia al momento de crear objetos de solicitudes, usuarios o servicios \cite{factory2018implementation}, permitiendo una instanciación más controlada y flexible. En sistemas orientados a la escalabilidad, como aquellos basados en Laravel o microservicios, este patrón reduce la dependencia directa del código con las clases concretas, favoreciendo la inyección de dependencias y la facilidad de pruebas unitarias.

La gestión de bases de datos constituye el núcleo sobre el cual se sostiene la lógica de persistencia de cualquier sistema \cite{mysql2023handbook}. Los fundamentos sobre la gestión de bases de datos no se limitan al almacenamiento de información, sino que abarcan la integridad referencial, la normalización, la optimización de consultas y la seguridad de los datos \cite{phpmyadmin2021security}. En este sentido, herramientas como MySQL y PHPMyAdmin continúan siendo referentes en la administración de datos estructurados, mientras que su integración con ORM (Object Relational Mapping) en frameworks como Laravel \cite{orm2022comparison} simplifica la manipulación de datos a través de modelos representativos de las entidades del dominio.